\documentclass[11pt,reqno]{amsart}

\usepackage[T1]{fontenc}
\usepackage[utf8]{inputenc}
\usepackage{lmodern}
\usepackage{microtype}
\usepackage{amssymb,amsmath,amsfonts,amsthm,mathtools,mathrsfs}
\usepackage{aliascnt}
\usepackage[dvipsnames]{xcolor}
\usepackage[margin=1in]{geometry}
\usepackage[shortlabels]{enumitem}
\usepackage{booktabs,tabularx,array}
\usepackage{tikz}
\usepackage[colorlinks=true]{hyperref}
\hypersetup{
  linkcolor=blue!45!black,
  citecolor=magenta,
  urlcolor=blue!45!black,
  pdftitle={An entropy bound for the attractive logarithmic gas},
  pdfauthor={Matias G. Delgadino, Rishabh Gvalani, Matthew Rosenzweig}
}
\usepackage[nameinlink,capitalise,noabbrev]{cleveref}

\raggedbottom
\allowdisplaybreaks
\numberwithin{equation}{section}
\emergencystretch=3em
\setlist{listparindent=\parindent,parsep=0pt}
\setlist[itemize]{topsep=4pt,itemsep=2pt,parsep=1pt,leftmargin=1.8em}
\setlist[enumerate]{topsep=4pt,itemsep=3pt,parsep=1pt,leftmargin=2.2em}

\theoremstyle{plain}
\newtheorem{thm}{Theorem}[section]
\newaliascnt{prop}{thm}
\newtheorem{prop}[prop]{Proposition}
\aliascntresetthe{prop}
\newaliascnt{lemma}{thm}
\newtheorem{lemma}[lemma]{Lemma}
\aliascntresetthe{lemma}
\newaliascnt{cor}{thm}
\newtheorem{cor}[cor]{Corollary}
\aliascntresetthe{cor}
\newaliascnt{assumption}{thm}
\newtheorem{assumption}[assumption]{Assumption}
\aliascntresetthe{assumption}
\theoremstyle{definition}
\newaliascnt{mydef}{thm}
\newtheorem{mydef}[mydef]{Definition}
\aliascntresetthe{mydef}
\newaliascnt{remark}{thm}
\newtheorem{remark}[remark]{Remark}
\aliascntresetthe{remark}
\newaliascnt{conj}{thm}
\newtheorem{conj}[conj]{Conjecture}
\aliascntresetthe{conj}

\Crefname{thm}{Theorem}{Theorems}
\Crefname{prop}{Proposition}{Propositions}
\Crefname{lemma}{Lemma}{Lemmas}
\Crefname{cor}{Corollary}{Corollaries}
\Crefname{assumption}{Assumption}{Assumptions}
\Crefname{mydef}{Definition}{Definitions}
\Crefname{remark}{Remark}{Remarks}
\Crefname{conj}{Conjecture}{Conjectures}
\crefformat{equation}{(#2#1#3)}

\DeclarePairedDelimiter{\pa}{\lparen}{\rparen}
\DeclarePairedDelimiter{\brak}{\lbrack}{\rbrack}
\DeclarePairedDelimiter{\brac}{\lbrace}{\rbrace}
\DeclarePairedDelimiter{\abs}{\lvert}{\rvert}
\DeclarePairedDelimiter{\norm}{\lVert}{\rVert}
\DeclarePairedDelimiter{\ip}{\langle}{\rangle}

\renewcommand{\d}{\mathsf d}
\newcommand{\dd}{\mathop{}\!\mathrm d}
\newcommand{\T}{\mathbb T}
\newcommand{\R}{\mathbb R}
\newcommand{\C}{\mathbb C}
\newcommand{\Z}{\mathbb Z}
\newcommand{\N}{\mathbb N}
\newcommand{\E}{\mathbb E}
\newcommand{\Pb}{\mathbb P}
\newcommand{\Q}{\mathbb Q}
\newcommand{\Id}{\operatorname{Id}}
\newcommand{\Tr}{\operatorname{Tr}}
\newcommand{\TV}{\operatorname{TV}}
\DeclareMathOperator{\sgn}{sgn}
\newcommand{\rank}{\operatorname{rank}}
\newcommand{\HS}{\mathrm{HS}}
\newcommand{\op}{\mathrm{op}}
\newcommand{\diag}{\Delta}
\newcommand{\one}{\mathbf 1}
\newcommand{\be}{\beta}
\newcommand{\al}{\alpha}
\newcommand{\vep}{\varepsilon}
\newcommand{\g}{\mathsf g}
\newcommand{\K}{\mathsf K}
\newcommand{\Fr}{\mathsf F}
\newcommand{\Hr}{\mathsf H}
\newcommand{\Er}{\mathsf E}
\newcommand{\Ec}{\mathcal E}
\newcommand{\Pc}{\mathcal P}
\newcommand{\Ent}{\operatorname{Ent}}
\newcommand{\Gcal}{\mathcal G}
\newcommand{\Ecal}{\mathcal E}
\newcommand{\pcspip}[2]{\left\langle #1,#2\right\rangle}
\newcommand{\tl}{\widetilde}
\newcommand{\f}{\mathsf f}
\newcommand{\wh}{\widehat}
\newcommand{\muunif}{\mu_{\mathrm{unif}}}
\newcommand{\gm}{\mathrm{gm}}
\newcommand{\pf}{\mathrm{pf}}
\newcommand{\spec}{\mathrm{s}}
\newcommand{\coll}{\mathrm{c}}
\newcommand{\tr}{\mathrm{tr}}
\newcommand{\XN}{X_N}
\newcommand{\Part}{\mathfrak P}
\DeclareMathOperator*{\esssup}{ess\,sup}
\DeclareMathOperator{\diver}{div}
\newcommand{\step}[1]{\par\medskip\noindent\textbf{#1}}


\title[An entropy bound for the attractive logarithmic gas]{An entropy bound for the attractive logarithmic gas}
\author{M. G. Delgadino, R. Gvalani, and M. Rosenzweig}
\date{\today}
\begin{document}
\begin{abstract}
Assuming a logarithmic Hardy--Littlewood--Sobolev inequality, we prove a bound on the attractive logarithmic partition function that is uniform in the number of particles.  A positive-definite cutoff separates a smooth bulk, which inherits the entropy inequality, from a singular interaction supported near the diagonal.  The latter is perturbative when each box contains few particles; the complementary event is exponentially small.  We give the required estimate for centered arrays of independent variables and include a cell-projection proof of the continuous-kernel bound in an appendix.
\end{abstract}
\maketitle

\section{Setting and partition bound}
\label{sec:setting}
Let $\d\ge2$ and let $\muunif$ be normalized Lebesgue measure on $\T^\d=\R^\d/\Z^\d$.  The mean-zero periodic logarithmic kernel is normalized by
\begin{equation}
 \wh\g(k)=\frac{c_\d}{(2\pi|k|)^\d}\quad(k\ne0),
 \qquad c_\d=2^{\d-1}\pi^{\d/2}\Gamma(\d/2),
 \qquad \g(x)=-\log|x|+O_\d(1)
 \quad(x\to0).
 \label{eq:kernel-normalization}
\end{equation}
Here, $\wh f(k)=\int f(x)e^{-2\pi ik\cdot x}\,\dd\muunif(x)$ and $\wh\g(0)=0$.
For independent $X_1,\ldots,X_N$ with law $\muunif$, set
\begin{equation}
 \mathsf W_N(K)\coloneqq\frac1N\sum_{i\ne j}K(X_i,X_j),
 \qquad
 Z_{N,\be}\coloneqq\E\big[e^{\frac\be2\mathsf W_N(\g)}\big],
 \label{eq:homogeneous-partition-function}
\end{equation}
where a function of one variable denotes the translation-invariant kernel $K(x-y)$.

Our only macroscopic assumption is the log-HLS inequality: for some $B\in(0,2\d]$,
\begin{equation}
 \Hr(\rho\muunif\mid\muunif)
 \coloneqq\int\rho\log\rho\,\dd\muunif
 \ge\frac B2\iint\g(x-y)(\rho(x)-1)(\rho(y)-1)
                  \,\dd\muunif(x)\dd\muunif(y)
 \label{eq:strict-subcritical-macroscopic-inequality}
\end{equation}
for every probability density of finite entropy.  No value or optimality of $B$ is needed.

\begin{thm}[Uniform partition bound]
\label{thm:principal}
Under \eqref{eq:strict-subcritical-macroscopic-inequality}, for every $0\le\be_0<B$ there is a finite constant $C=C(\d,B,\be_0)$ such that
\begin{equation}
 \sup_{N\ge2}\sup_{0\le\be\le\be_0}Z_{N,\be}\le C.
 \label{eq:principal-order-one}
\end{equation}
\end{thm}

The proof splits $\g$ into a continuous bulk, controlled by Lemma~\ref{lem:continuous-entropy-partition}, and a localized singular part.  For the latter, the number of particles in a box makes the effective coupling small enough for Lemma~\ref{lem:canonical-array-exponential}.

\textbf{Boxes and the good event.}
For an integer $J\ge3$, partition the torus into
\begin{equation}
 A_a=J^{-1}(a+[0,1)^\d),\qquad
 a\in(\Z/J\Z)^\d,\qquad M=J^\d.
 \label{eq:common-box-partition}
\end{equation}
Write $L_i=a$ when $X_i\in A_a$, and $n_a=\#\{i:L_i=a\}$.
Let $\nu_a=M\one_{A_a}\muunif$ be the uniform probability measure in $A_a$, and let $Q_J$ denote box averaging: $(Q_Ju)|_{A_a}=\int u\,\dd\nu_a$.
These boxes and averages are used in both parts of the proof.

For the singular part, fix $0<\delta<1$ and choose the cutoff and the occupancy event together:
\begin{equation}
 r=\frac1{2J},\qquad
 \mathcal A_{N,J,\delta}=\{n_a\le\delta N\ \text{for every }a\},
 \qquad
 \mathcal B_{N,J,\delta}=\mathcal A_{N,J,\delta}^{\mathsf c}.
 \label{eq:physical-neighbor-events}
\end{equation}
The singular kernel constructed below is supported at distance at most $r$.
Nonadjacent boxes are separated by at least $J^{-1}=2r$, so only the same box or adjacent boxes can interact.  Here, adjacency includes shared faces, edges, and corners, with periodic labels: the allowed displacements are
$\mathcal V=\{-1,0,1\}^\d$.
On $\mathcal A_{N,J,\delta}$, one box contains at most $\delta N$ particles and two boxes at most $2\delta N$; see Figure~\ref{fig:periodic-cell-pair-band}.

The bulk lemma uses the same family of box averages, with a finer grid if needed after $r$ has been fixed.  All grids and cutoffs are fixed independently of $N$.

\begin{figure}[htbp]
\centering
\begin{tikzpicture}[x=.9cm,y=.9cm,font=\small]
\node at (7.2,5.25) {Periodic cell pairs on $\mathbb T\times\mathbb T$ (one-dimensional schematic)};
\definecolor{cellself}{RGB}{35,67,104}
\definecolor{cellnear}{RGB}{237,181,91}
\definecolor{cellfar}{RGB}{221,237,243}
\foreach \cells/\shift/\cutoff in {12/0/{1/24},24/5.2/{1/48},48/10.4/{1/96}} {
 \begin{scope}[shift={(\shift,0)}]
  \node[font=\bfseries] at (2,4.6) {$J=\cells$, $r=\cutoff$};
  \fill[cellfar] (0,0) rectangle (4,4);
  \pgfmathtruncatemacro{\lastcell}{\cells-1}
  \foreach \a in {0,...,\lastcell} {
   \foreach \offset in {-1,1} {
    \pgfmathtruncatemacro{\b}{mod(\a+\offset+\cells,\cells)}
    \fill[cellnear] ({4*\a/\cells},{4*\b/\cells}) rectangle ({4*(\a+1)/\cells},{4*(\b+1)/\cells});
   }
   \fill[cellself] ({4*\a/\cells},{4*\a/\cells}) rectangle ({4*(\a+1)/\cells},{4*(\a+1)/\cells});
  }
  \foreach \a in {0,...,\cells} {
   \draw[white,line width=.13pt] ({4*\a/\cells},0)--({4*\a/\cells},4);
   \draw[white,line width=.13pt] (0,{4*\a/\cells})--(4,{4*\a/\cells});
  }
  \draw[black!65,line width=.5pt] (0,0) rectangle (4,4);
  \node[below left,font=\scriptsize] at (0,0) {$0$};
  \node[below,font=\scriptsize] at (4,0) {$1$};
  \node[left,font=\scriptsize] at (0,4) {$1$};
  \node[below] at (2,-.08) {$x\in\mathbb T$};
  \ifnum\cells=12\node[rotate=90] at (-.42,2) {$y\in\mathbb T$};\fi
  \node at (2,-.85) {$J^{-1}=2r$};
 \end{scope}
}
\foreach \shift in {0,5.2} {\draw[->,thick,black!65] ({4.25+\shift},2)--({4.95+\shift},2);}
\fill[cellself] (0,-1.6) rectangle (.24,-1.36);
\node[anchor=west] at (.4,-1.48) {Same box: at most $\delta N$ particles};
\fill[cellnear] (0,-2.12) rectangle (.24,-1.88);
\node[anchor=west] at (.4,-2) {Two adjacent boxes: at most $2\delta N$ particles};
\fill[cellfar,draw=black!25] (0,-2.64) rectangle (.24,-2.4);
\node[anchor=west] at (.4,-2.52) {Nonadjacent boxes: the raw singular kernel $s_r$ vanishes};
\end{tikzpicture}

\caption{The grid is chosen to match the singular cutoff: the cell side is $2r$.  In this one-dimensional schematic, $s_r(x-y)$ can be nonzero only in a diagonal square or one of its two neighboring squares, including the periodic corner squares.  The colored boxes cover its support; they are not all filled by it.  On $\mathcal A_{N,J,\delta}$, the one- and two-box counts are bounded by $\delta N$ and $2\delta N$.  Refining the grid and reducing $r$ together narrows the physical band while keeping three possible partners per box ($3^\d$ in dimension $\d$).}
\label{fig:periodic-cell-pair-band}
\end{figure}

\section{Separate the bulk and the diagonal}
\label{sec:physical-cutoff-proof}
With $r=(2J)^{-1}$ as above, choose a nonzero, nonnegative, radial $\chi\in C_c^\infty(B_{1/2})$ and set $\phi=\chi*\chi$.  Define
\begin{equation}
 \begin{aligned}
 s_r(x)&=\frac1{\phi(0)}\sum_{m\in\Z^\d}\int_0^r
                  \phi\!\left(\frac{x+m}{t}\right)\frac{\dd t}{t},\\
 c_r&=\int s_r\,\dd\muunif,\qquad
 \widetilde s_r=s_r-c_r,\qquad b_r=\g-\widetilde s_r.
 \end{aligned}
 \label{eq:compact-positive-log-split}
\end{equation}
The raw kernel $s_r$ is nonnegative, supported in $\{\operatorname{dist}_{\T^\d}(x,0)\le r\}$, and $c_r=O(r^\d)$.  Since $\wh\phi=|\wh\chi|^2\ge0$, radial Fourier inversion gives, for $k\ne0$,
\begin{equation}
 0\le\wh{\widetilde s_r}(k)
 =\frac1{\phi(0)}\int_0^r t^{\d-1}\wh\phi(tk)\,\dd t
 \le\frac1{\phi(0)}\int_0^\infty t^{\d-1}\wh\phi(tk)\,\dd t
 =\wh\g(k).
 \label{eq:compact-split-fourier-order}
\end{equation}
Thus both centered parts are positive semidefinite.  The Fourier coefficients of $b_r$, given by the integral from $r$ to infinity, decay rapidly, so $b_r$ is smooth.  Its energy is bounded by that of $\g$; hence it inherits \eqref{eq:strict-subcritical-macroscopic-inequality} with the same $B$.

Group the label vectors $L$ by their occupancies $n=(n_a)$.  For a fixed $n$, retag the particles as $Y_{a,1},\ldots,Y_{a,n_a}$ in box $A_a$, and let $\E_n$ denote expectation under their independent product law, with each $Y_{a,p}$ distributed according to $\nu_a$.  Since the energy is symmetric in the particles,
\begin{equation}
 Z_{N,\be}
 =\sum_{\substack{n\in\N_0^M\\\sum_a n_a=N}}
   \pi_N(n)\,\E_n\big[e^{\frac\be2\mathsf W_N(\g)}\big],
 \qquad
 \pi_N(n)=\frac{N!}{M^N\prod_a n_a!}.
 \label{eq:partition-occupancy-sum}
\end{equation}
Here, $\mathsf W_N$ is evaluated on the retagged particles in any order, and the multinomial weights $\pi_N(n)$ sum to one.  Split this sum into good vectors, with $\max_a n_a\le\delta N$, and the remaining bad vectors.  Fix $0<\be_0<B$ and conjugate $p,q>1$ with $p\be_0<B$.  H\"older on the good contribution gives
\begin{equation}
 \begin{aligned}
 Z_{N,\be}
 &\le
 \Big(\E\big[e^{\frac{p\be}{2}\mathsf W_N(b_r)}\big]\Big)^{1/p}
 \Big(\E\big[\one_{\mathcal A_{N,J,\delta}}
        e^{\frac{q\be}{2}\mathsf W_N(\widetilde s_r)}\big]\Big)^{1/q}\\
 &\qquad+\E\big[\one_{\mathcal B_{N,J,\delta}}
                     e^{\frac\be2\mathsf W_N(\g)}\big],
 \qquad 0\le\be\le\be_0.
 \end{aligned}
 \label{eq:physical-cutoff-holder}
\end{equation}
The indicator stays on the singular factor, whose coupling $q\be$ can be large.  The next three sections bound the bad event, the bulk, and the singular factor.

\section{Discard overcrowded boxes}
\label{subsec:proof-bad-set}
We first need a rough bound on the full Gibbs weight.  For $0<b<2\d$, let
$A_b=\int e^{b\g/2}\,\dd\muunif<\infty$.
Convexity of the exponential, followed by integration conditional on $X_i$, gives
\begin{equation}
 \begin{aligned}
 Z_{N,b}
 &=\E\left[\exp\left\{\frac1N\sum_{i=1}^N
                   \frac b2\sum_{j\ne i}\g(X_i-X_j)\right\}\right]\\
 &\le\frac1N\sum_{i=1}^N
       \E\left[e^{\frac b2\sum_{j\ne i}\g(X_i-X_j)}\right]
 =A_b^{N-1}.
 \end{aligned}
 \label{eq:rough-partition-bound}
\end{equation}
This uses only the local logarithmic singularity.  Moreover, $A_b\ge1$ by Jensen, and H\"older gives $Z_{N,\gamma}\le Z_{N,b}^{\gamma/b}\le A_b^{N-1}$ for $0\le\gamma\le b$.

Each $n_a$ is binomial with parameters $N,M^{-1}$.  For $M^{-1}<\delta$, Chernoff's inequality and a union bound give
\begin{equation}
 \begin{aligned}
 \Pb(\mathcal B_{N,J,\delta})
 &\le M e^{-N\mathsf D(\delta\Vert M^{-1})},\\
 \mathsf D(\delta\Vert M^{-1})
 &=\delta\log(\delta M)
   +(1-\delta)\log\frac{1-\delta}{1-M^{-1}}
 \longrightarrow\infty\quad(J\to\infty).
 \end{aligned}
 \label{eq:physical-bad-probability}
\end{equation}
Choose $s>1$ with $b=s\be_0<2\d$ and $s'=s/(s-1)$.  H\"older and \eqref{eq:rough-partition-bound} yield
\begin{equation}
 \E\big[\one_{\mathcal B_{N,J,\delta}}e^{\frac\be2\mathsf W_N(\g)}\big]
 \le M^{1/s'}\exp\left\{N
 \left(\frac{\log A_b}{s}
       -\frac{\mathsf D(\delta\Vert M^{-1})}{s'}\right)\right\}.
 \label{eq:physical-bad-gibbs-bound}
\end{equation}
For fixed $\delta>0$, choose $J$ large enough to make the exponent negative, then keep this grid fixed as $N$ varies.

\section{The continuous bulk}
\label{sec:continuous-localized-reduction}
The next lemma records a continuous-kernel entropy--energy partition bound; see Ren and Wang \cite[Corollary~4 and Remark~3]{RenWang2025} for the bounded-interaction torus version.  For completeness, Appendix~\ref{app:continuous-proof} gives a proof using cell averages, an energy projection, and the estimate of Section~\ref{sec:canonical-arrays}.

\begin{lemma}[Continuous kernels controlled by entropy]
\label{lem:continuous-entropy-partition}
Let $K\in C(\T^\d\times\T^\d)$ be real and symmetric, with $\int K(x,y)\,\dd\muunif(y)=0$.  Suppose its integral operator $T_K$ is nonnegative and, for some $B>0$,
\begin{equation}
 \Hr(\rho\muunif\mid\muunif)
 \ge\frac B2\ip{\rho-1,T_K(\rho-1)}
 \quad\text{for every finite-entropy probability density }\rho.
 \label{eq:continuous-entropy-assumption}
\end{equation}
Then, for every $0\le\be_0<B$,
\begin{equation}
 \sup_{N\ge1}\sup_{0\le\be\le\be_0}
 \E\big[e^{\frac\be2\mathsf W_N(K)}\big]<\infty.
 \label{eq:continuous-partition-conclusion}
\end{equation}
\end{lemma}

For the comparison, set $W_+=0$, $W_-=\be K$, $m_*=\muunif$, and $\lambda=0$ in their entropy hypothesis, with $0<\varepsilon\le1/2$ and $(1+\varepsilon)\be<B$.  The diagonal included there contributes at most $\be\|K\|_\infty/2$ to the exponent.

For $K(x,y)=b_r(x-y)$, the hypotheses follow from \eqref{eq:compact-split-fourier-order} and log-HLS.  In particular,
\begin{equation}
 \sup_{N\ge1}\sup_{0\le\be\le\be_0}
 \E\big[e^{\frac{p\be}{2}\mathsf W_N(b_r)}\big]<\infty
 \qquad\text{for every fixed }0<r<1/4.
 \label{eq:physical-continuous-bulk-bound}
\end{equation}

\section{The singular part and the conclusion}
\label{sec:local-singular}
The good event makes the singular factor perturbative through the number of particles in each interacting pair of boxes.  The moment estimate used below is proved in Lemma~\ref{lem:canonical-array-exponential}.
\begin{lemma}[Short-range estimate under a local count bound]
\label{lem:physical-local-count-bound}
For every $0<t_0<\infty$, there are $0<\delta_0<1$ and an integer $J_0\ge3$ such that, for each fixed $0<\delta\le\delta_0$ and integer $J\ge J_0$, with $r=(2J)^{-1}$,
\begin{equation}
 \sup_{N\ge1}\sup_{0\le t\le t_0}
 \E\Big[\one_{\mathcal A_{N,J,\delta}}
 e^{\frac t2\mathsf W_N(\widetilde s_r)}\Big]<\infty.
 \label{eq:physical-local-count-conclusion}
\end{equation}
The constant may depend on $J,\delta,t_0$; these parameters are independent of $N$.
\end{lemma}

\begin{proof}
Fix $0\le t\le t_0$.  We separate off the box averages, then use \eqref{eq:partition-occupancy-sum} to estimate the remaining integral for each good occupancy vector.

\textbf{1. Split the kernel and apply H\"older.}
Apply the box average $Q_J$ in each variable and write
\begin{equation}
 \widetilde s_r=a_r+h_r,
 \qquad h_r=(I-Q_{J,x})(I-Q_{J,y})s_r,
 \qquad a_r=Q_{J,x}s_r+Q_{J,y}s_r-Q_{J,x}Q_{J,y}s_r-c_r.
 \label{eq:physical-local-centering}
\end{equation}
Here, $h_r$ has zero average in each variable within each box.  Both $a_r$ and $h_r$ also have zero average in each variable under $\muunif$.  The particle sum splits exactly as
\begin{equation}
 \frac t2\mathsf W_N(\widetilde s_r)
 =\frac tN\sum_{i<j}a_r(X_i,X_j)
  +\frac tN\sum_{i<j}h_r(X_i,X_j).
 \label{eq:physical-particle-sum-splitting}
\end{equation}
H\"older with exponents $2,2$ therefore gives
\begin{equation}
 \E\big[\one_{\mathcal A_{N,J,\delta}}e^{\frac t2\mathsf W_N(\widetilde s_r)}\big]
 \le
 \Big(\E\big[e^{t\mathsf W_N(a_r)}\big]\Big)^{1/2}
 \Big(\E\big[\one_{\mathcal A_{N,J,\delta}}e^{t\mathsf W_N(h_r)}\big]\Big)^{1/2}.
 \label{eq:physical-local-holder}
\end{equation}
We bound the first factor by taking $r$ small and the second by taking $\delta$ small. 

\textbf{2. The averaging term is perturbative.}
Every box average of the nonnegative kernel $s_r$ is at most $Mc_r$.  Since $c_r\le Cr^\d$ and $Mr^\d=2^{-\d}$, we have
\begin{equation}
 \|a_r\|_\infty\le C,
 \qquad
 \sup_x\int|a_r(x,y)|\,\dd\muunif(y)\le4c_r\le Cr^\d.
 \label{eq:physical-average-smallness}
\end{equation}
The second estimate follows because each averaging term has integral $c_r$ in either variable.  These bounds imply
$\sup_x\int(e^{|a_r(x,y)|/\kappa}-1)\,\dd\muunif(y)
\le Cr^\d(e^{C/\kappa}-1)$.
Taking $\kappa(r)=C'/\log(1/r)$ with $C'$ large enough makes the right side at most one for small $r$.  Expanding the exponential and using symmetry gives, for every integer $p\ge1$,
\begin{equation}
 \sup_x\int|a_r(x,y)|^p\,\dd\muunif(y)
 +\sup_y\int|a_r(x,y)|^p\,\dd\muunif(x)
 \le2\kappa(r)^p p!,
 \qquad \kappa(r)\longrightarrow0.
 \label{eq:physical-average-factorial}
\end{equation}
Apply Lemma~\ref{lem:canonical-array-exponential} with $m=N$, $\kappa=\kappa(r)$, and $\theta=2t$:
\begin{equation}
 \E\big[e^{t\mathsf W_N(a_r)}\big]
 =\E\Big[e^{\frac{2t}{N}\sum_{i<j}a_r(X_i,X_j)}\Big]
 \le\frac{2}{1-(2tC_{\mathrm u}\kappa(r))^2}
 \le\frac83,
 \qquad 0\le t\le t_0,
 \label{eq:physical-average-exponential-bound}
\end{equation}
provided $J\ge J_0$, with $J_0$ chosen so that $2t_0C_{\mathrm u}\kappa((2J_0)^{-1})\le1/2$.

\textbf{3. Few particles make the singular term perturbative.}
Under the box-uniform measures $\nu_a$, there is a constant $\kappa_*$, independent of $J$, such that, for every integer $p\ge1$,
\begin{equation}
 \sup_{y\in A_b}\int_{A_a}|h_r(x,y)|^p\,\dd\nu_a(x)
 +\sup_{x\in A_a}\int_{A_b}|h_r(x,y)|^p\,\dd\nu_b(y)
 \le2\kappa_*^p p!,
 \qquad p\ge1.
 \label{eq:physical-singular-factorial}
\end{equation}
Indeed, Cauchy--Schwarz gives $0\le\phi\le\phi(0)$ in \eqref{eq:compact-positive-log-split}, so $s_r(z)\le(\log(r/\operatorname{dist}_{\T^\d}(z,0)))_+$.  Thus
$M\int s_r^p\,\dd\muunif\le M|B_r|p!/\d^p\le C^p p!$.
The averages subtracted to form $h_r$ are bounded by $Mc_r\le C$, which proves the same estimate for $h_r$ with a larger constant.

The occupancy decomposition gives exactly
\begin{equation}
 \E\big[\one_{\mathcal A_{N,J,\delta}}e^{t\mathsf W_N(h_r)}\big]
 =\sum_{\substack{n\in\N_0^M,\ \sum_a n_a=N\\\max_a n_a\le\delta N}}
   \pi_N(n)\,\E_n\big[e^{t\mathsf W_N(h_r)}\big].
 \label{eq:physical-good-occupancy-sum}
\end{equation}
Fix one vector in this sum.  Its particles form $M$ independent groups, each of size at most $\delta N$.  Only the same or adjacent boxes interact: $s_r$ vanishes on other box products, and box averaging preserves this property for $h_r$.  Thus each group interacts with at most $K-1$ others, where $K=3^\d$.

Apply the grouped conclusion \eqref{eq:canonical-grouped-exponential} of Lemma~\ref{lem:canonical-array-exponential}, with $m=N$, $m_*=\max_a n_a$, $\kappa=\kappa_*$, and $u=2t$ since $\mathsf W_N(h_r)=2N^{-1}\sum_{i<j}h_r(X_i,X_j)$.  Choose $\delta_0$ so that $8K C_{\mathrm u}\kappa_*t_0\delta_0\le1/2$.  Then
\begin{equation}
 4K C_{\mathrm u}\kappa_*(2t)\frac{m_*}{N}
 \le8K C_{\mathrm u}\kappa_*t_0\delta\le\frac12,
 \qquad
 \E_n\big[e^{t\mathsf W_N(h_r)}\big]\le\left(\frac83\right)^M.
 \label{eq:physical-local-effective-coupling}
\end{equation}
This is uniform over all good occupancy vectors.  Their multinomial weights sum to at most one, so \eqref{eq:physical-good-occupancy-sum} yields
\begin{equation}
 \sup_{N\ge1}\sup_{0\le t\le t_0}
 \E\big[\one_{\mathcal A_{N,J,\delta}}e^{t\mathsf W_N(h_r)}\big]
 \le\left(\frac83\right)^M<\infty.
 \label{eq:physical-centered-local-bound}
\end{equation}
Finally, insert this estimate and \eqref{eq:physical-average-exponential-bound} into \eqref{eq:physical-local-holder}.  The result is at most $(8/3)^{(M+1)/2}$, uniformly in $N\ge1$ and $0\le t\le t_0$.  Both $J_0$ and $\delta_0$ are chosen independently of $N$.
\end{proof}

\begin{proof}[Conclusion of Theorem~\ref{thm:principal}]
The case $\be_0=0$ is immediate.  Otherwise, choose $p,q$ as in \eqref{eq:physical-cutoff-holder}, and take $\delta\le\delta_0$ from Lemma~\ref{lem:physical-local-count-bound} with $t_0=q\be_0$.  Next, choose $J\ge J_0$ large enough that $J^{-\d}<\delta$ and \eqref{eq:physical-bad-gibbs-bound} decays exponentially, and set $r=(2J)^{-1}$.  The singular factor is bounded by Lemma~\ref{lem:physical-local-count-bound}, and the bulk factor by \eqref{eq:physical-continuous-bulk-bound}.  Inserting these bounds into \eqref{eq:physical-cutoff-holder} proves \eqref{eq:principal-order-one}.
\end{proof}

\section{A perturbative estimate for centered arrays}
\label{sec:canonical-arrays}
The variables in a local block have different box-uniform laws.  The following estimate requires independence and centering, but not identical distributions.

\begin{lemma}[Canonical-array exponential estimate]
\label{lem:canonical-array-exponential}
There is a universal $C_{\mathrm u}<\infty$ with the following property.  Let $m\ge1$, let $Y_1,\ldots,Y_m$ be independent, and let $h_{ij}$, $i<j$, be real measurable kernels centered in both variables:
\begin{equation}
 \E[h_{ij}(Y_i,y)]=0,\qquad \E[h_{ij}(x,Y_j)]=0
 \quad\text{for almost every }y,x.
 \label{eq:canonical-array-centering}
\end{equation}
Suppose, for some $\kappa>0$, every integer $p\ge1$, and every $i<j$,
\begin{equation}
 \esssup_y\E[|h_{ij}(Y_i,y)|^p]
 +\esssup_x\E[|h_{ij}(x,Y_j)|^p]\le2\kappa^p p!.
 \label{eq:canonical-array-kernel-moments}
\end{equation}
For $S=\sum_{i<j}h_{ij}(Y_i,Y_j)$, every integer $k\ge1$, and real $\theta$, one has
\begin{equation}
 \E[|S|^{2k}]\le(C_{\mathrm u}\kappa m)^{2k}(2k)!,
 \qquad
 \E[e^{\theta S/m}]
 \le\frac2{1-(C_{\mathrm u}\kappa|\theta|)^2}
 \quad\text{if }C_{\mathrm u}\kappa|\theta|<1.
 \label{eq:canonical-array-exponential}
\end{equation}
Zero kernels are allowed.  We will use the following grouped form.  Partition the $m$ variables into $M$ groups, each containing at most $m_*$ variables.  Suppose each group interacts with at most $K-1$ other groups, where $K\ge1$ is an integer: kernels between all remaining pairs of groups vanish.  Then, for $N\ge1$ and $u\in\R$,
\begin{equation}
 \E[e^{uS/N}]\le\left(\frac83\right)^M
 \quad\text{if}\quad
 4K C_{\mathrm u}\kappa|u|\frac{m_*}{N}\le\frac12.
 \label{eq:canonical-grouped-exponential}
\end{equation}
\end{lemma}

\begin{proof}
We give the counting argument.  Extend the kernels by $h_{ji}(y,x)=h_{ij}(x,y)$.  Independently of the variables $Y_i$, place each index in $I$ independently with probability $1/2$, and put
$D_I=\sum_{i\in I,j\notin I}h_{ij}(Y_i,Y_j)$.
Keep the $Y_i$ fixed and average only over the choice of $I$, denoting this average by $\E_I$.  The transpose convention gives
\begin{equation}
 \begin{aligned}
 D_I&=\sum_{i<j}
 \big(\one_{\{i\in I,\,j\notin I\}}+\one_{\{j\in I,\,i\notin I\}}\big)
 h_{ij}(Y_i,Y_j),\\
 \E_I[D_I]&=\sum_{i<j}\left(\frac14+\frac14\right)h_{ij}(Y_i,Y_j)
 =\frac12 S.
 \end{aligned}
 \label{eq:array-split-identity}
\end{equation}
Each of the two disjoint ways to separate a pair has probability $(1/2)(1/2)=1/4$.  Thus $S=2\E_I[D_I]$ for every fixed sample.  Jensen's inequality, followed by expectation over the $Y_i$, now gives
\begin{equation}
 \E[|S|^{2k}]\le2^{2k}\E_I\E[|D_I|^{2k}].
 \label{eq:array-random-split}
\end{equation}
Fix $I$ and expand $D_I^{2k}$.  Each term is indexed by a row word
$\mathbf i=(i_1,\ldots,i_{2k})$ and a column word
$\mathbf j=(j_1,\ldots,j_{2k})$.
If any label occurs only once in either word, centering makes the expectation zero.  A surviving word therefore uses at most $k$ labels.

For each label $\ell\in\{1,\ldots,m\}$, let
\begin{equation}
 a_\ell=\#\{v\in\{1,\ldots,2k\}:i_v=\ell\}
 \label{eq:row-multiplicity}
\end{equation}
be the multiplicity of $\ell$ in the row word $\mathbf i$, so that $\sum_\ell a_\ell=2k$.  Since every $a_\ell\ne1$ for a surviving word, the labels with $a_\ell\ge2$ are the active ones, and there are at most $k$ of them; factors with $a_\ell=0$ equal $1$ in all products and sums over $\ell$ below.  Condition on all column variables.  Since every row index lies in $I$ and every column index lies outside $I$, each factor becomes a function of a single row variable, and the factors sharing the row label $\ell$ form a group of $a_\ell$ functions of $Y_\ell$.  By independence of the row variables, the conditional expectation factors over these groups.  Within the group of label $\ell$, H\"older with all exponents equal to $a_\ell$ and the moment hypothesis \eqref{eq:canonical-array-kernel-moments} bound the group by $2\kappa^{a_\ell}a_\ell!$, uniformly in the frozen column values.  In formulas, write $T=\prod_{v=1}^{2k}h_{i_vj_v}(Y_{i_v},Y_{j_v})$ and let $\mathcal C=\sigma(Y_j:j\notin I)$.  The tower property gives $\E[T]=\E\big[\E[T\mid\mathcal C]\big]$.  The row variables are independent of $\mathcal C$ and of one another, so
\begin{equation}
 \E[T\mid\mathcal C]
 =\prod_\ell\Phi_\ell(Y_{j_v}:j_v\notin I),
 \qquad
 \Phi_\ell(y)=\E\Big[\prod_{v:\,i_v=\ell}h_{\ell j_v}(Y_\ell,y_{j_v})\Big].
 \label{eq:array-conditional-factor}
\end{equation}
Hölder with exponents $a_\ell$ and \eqref{eq:canonical-array-kernel-moments} give, for almost every $y$,
\begin{equation}
 |\Phi_\ell(y)|
 \le\prod_{v:\,i_v=\ell}\E\big[|h_{\ell j_v}(Y_\ell,y_{j_v})|^{a_\ell}\big]^{1/a_\ell}
 \le\big(2\kappa^{a_\ell}a_\ell!\big)^{a_\ell/a_\ell}
 =2\kappa^{a_\ell}a_\ell!.
 \label{eq:array-holder-group}
\end{equation}
Since at most $k$ labels survive, $|\E[T]|\le\E\big[|\E[T\mid\mathcal C]|\big]\le\prod_\ell2\kappa^{a_\ell}a_\ell!$, that is,
\begin{equation}
 \left|\E\left[\prod_{v=1}^{2k}
 h_{i_vj_v}(Y_{i_v},Y_{j_v})\right]\right|
 \le 2^k\kappa^{2k}\prod_\ell a_\ell!.
 \label{eq:array-word-moment}
\end{equation}
There are $(2k)!/\prod_\ell a_\ell!$ row words with these multiplicities, so their factorials cancel.  Write $A$ for the number of surviving multiplicity vectors on at most $m$ labels, and $B_k$ for the number of surviving words.  Summing absolute values yields
\begin{equation}
 \E[|D_I|^{2k}]\le2^k\kappa^{2k}(2k)!\,A B_k.
 \label{eq:array-count-reduction}
\end{equation}

It remains to bound $A B_k$.  If $m<2k$, counting all weak compositions gives
$A\le\binom{2k+m-1}{m-1}\le2^{4k}$, and $B_k\le m^{2k}$.
If $m\ge2k$, choose the $r\le k$ active labels.  There are at most $4^k$ possible positive multiplicity vectors on these labels and at most $k^{2k}$ words.  Since
$\sum_{r=0}^k\binom mr\le(em/k)^k$, we obtain
\begin{equation}
 A\le4^k(em/k)^k,\qquad
 B_k\le k^{2k}(em/k)^k,\qquad
 A B_k\le(4e^2)^k m^{2k}.
 \label{eq:array-combinatorial-bound}
\end{equation}
Both cases, together with \eqref{eq:array-random-split}--\eqref{eq:array-count-reduction}, prove the moment bound after enlarging a universal constant.

Finally, $e^x\le2\cosh x$ and the even-moment bound imply
\begin{equation}
 \E[e^{\theta S/m}]
 \le2\sum_{k=0}^\infty
 \frac{|\theta|^{2k}\E[|S|^{2k}]}{m^{2k}(2k)!}
 \le2\sum_{k=0}^\infty(C_{\mathrm u}\kappa|\theta|)^{2k}.
 \label{eq:array-exponential-series}
\end{equation}
Tonelli justifies the nonnegative series.

For the grouped conclusion, draw an edge between each interacting pair of groups.  The graph has degree at most $K-1$.  Greedily color its edges with at most $2K-1$ colors so that edges sharing a group have different colors; place all within-group terms in one further class.  This writes $S=\sum_{c=1}^{2K}\sum_b S_{c,b}$, allowing empty classes.  Within each class, the sums $S_{c,b}$ use disjoint groups and are independent.  Each uses at most $2m_*$ variables.  Apply the one-block estimate with coefficient $2Ku/N$, then H\"older over the $2K$ classes:
\begin{equation}
 \E[e^{uS/N}]
 \le\prod_{c=1}^{2K}
      \left(\prod_b\E[e^{2KuS_{c,b}/N}]\right)^{1/(2K)}
 \le\left(\frac83\right)^M.
 \label{eq:array-grouped-holder}
\end{equation}
Indeed, the effective coupling in each factor is at most $4K|u|m_*/N$, and each class has at most $M$ nonempty blocks.  This proves \eqref{eq:canonical-grouped-exponential}.
\end{proof}

\appendix
\section{A cell-projection proof of the continuous-kernel bound}
\label{app:continuous-proof}
We prove Lemma~\ref{lem:continuous-entropy-partition}.  The steps are to replace a sample by its histogram and within-cell fluctuation, complete the square in the energy, and estimate the two resulting terms.  The histogram gives a finite-dimensional Gaussian integral; continuity makes the remaining centered kernel small.

\textbf{1. Histogram and energy projection.}
Use the partitions and box averages of Section~\ref{sec:setting}, choosing the resolution $J$ here for the fixed kernel $K$.  In the application $K=b_r$, the cutoff $r$ stays fixed while this auxiliary grid is refined.  With the corresponding labels and counts, set
\begin{equation}
 \mu_N=\frac1N\sum_i\delta_{X_i},\qquad
 f_L=\frac MN\sum_a n_a\one_{A_a},\qquad
 \nu_N=\mu_N-f_L\muunif.
 \label{eq:empirical-histogram-splitting}
\end{equation}
Thus $\nu_N(A_a)=0$.  Figure~\ref{fig:occupancy-and-offsets} shows this decomposition.

\begin{figure}[htbp]
\centering
\begin{tikzpicture}[x=1.12cm,y=1.12cm,font=\scriptsize,
  particle/.style={circle,inner sep=1.5pt},
  every path/.style={line width=.45pt}]
 \definecolor{cella}{RGB}{39,104,157}
 \definecolor{cellb}{RGB}{191,111,35}
 \definecolor{cellc}{RGB}{48,135,116}
 \definecolor{celld}{RGB}{133,93,158}
 \node at (1.1,2.95) {Sample from $\muunif^{\otimes N}$};
 \node at (5.1,2.95) {Graph of $f_L$};
 \node at (9.35,2.95) {Positions within cells};
 % Original sample: four cells, eight labelled particles.
 \foreach \x/\y/\col in {0/0/cella,1.1/0/cellb,0/1.1/cellc,1.1/1.1/celld}
   \fill[\col!6] (\x,\y) rectangle +(1.1,1.1);
 \draw (0,0) rectangle (2.2,2.2);
 \draw[gray] (1.1,0)--(1.1,2.2) (0,1.1)--(2.2,1.1);
 \foreach \x/\y/\col/\id in {.242/.242/cella/1,.770/.385/cella/4,.473/.792/cella/7,1.705/.528/cellb/2,.275/1.485/cellc/3,.748/1.837/cellc/8,1.342/1.848/celld/5,1.826/1.386/celld/6} {
   \node[particle,fill=\col] at (\x,\y) {};
   \node[anchor=west,inner sep=2pt] at (\x,\y) {$\id$};
 }
 \node at (1.1,-.35) {$x_i\in A_a$, cell side $J^{-1}$};
 \draw[->] (2.5,1.1)--(3.6,1.1);
 \node[align=center] at (3.05,1.65) {group\\by cell};
 % Raised histogram: horizontal coordinates label cells, height is density.
 \begin{scope}[shift={(4.02,.20)},x={(.80cm,-.13cm)},
               y={(.37cm,.80cm)},z={(0cm,.92cm)}]
  % Draw the back row first; small gaps keep all four columns legible.
  \foreach \xx/\yy/\hh/\col in {0/1/1/cellc,1/1/1/celld,0/0/1.5/cella,1/0/.5/cellb} {
   \filldraw[fill=\col!38,draw=\col!75!black]
     (\xx+.04,\yy+.04,0)--(\xx+.96,\yy+.04,0)--
     (\xx+.96,\yy+.04,\hh)--(\xx+.04,\yy+.04,\hh)--cycle;
   \filldraw[fill=\col!55,draw=\col!75!black]
     (\xx+.96,\yy+.04,0)--(\xx+.96,\yy+.96,0)--
     (\xx+.96,\yy+.96,\hh)--(\xx+.96,\yy+.04,\hh)--cycle;
   \filldraw[fill=\col!15,draw=\col!75!black]
     (\xx+.04,\yy+.04,\hh)--(\xx+.96,\yy+.04,\hh)--
     (\xx+.96,\yy+.96,\hh)--(\xx+.04,\yy+.96,\hh)--cycle;
  }
  \node at (.5,1.82,1) {$1$};
  \node at (1.5,1.5,1) {$1$};
  \node at (.5,.5,1.5) {$3/2$};
  \node at (1.5,.5,.5) {$1/2$};
 \end{scope}
 \node at (5.1,-.35) {Height $=Mn_a/N$};
 \node[font=\large] at (7.05,1.1) {$+$};
 % Same points in four separate copies of the unit square.
 \foreach \x/\y/\col/\tag in {8/0/cella/a,9.6/0/cellb/b,8/1.45/cellc/c,9.6/1.45/celld/d} {
   \fill[\col!6] (\x,\y) rectangle +(1.1,1.1);
   \draw[\col] (\x,\y) rectangle +(1.1,1.1);
   \node[anchor=south east,inner sep=2pt,text=\col] at (\x+1.1,\y) {$\tag$};
 }
 \foreach \x/\y/\col/\id in {8.242/.242/cella/1,8.770/.385/cella/2,8.473/.792/cella/3,10.205/.528/cellb/1,8.275/1.835/cellc/1,8.748/2.187/cellc/2,9.842/2.198/celld/1,10.326/1.736/celld/2} {
   \node[particle,fill=\col] at (\x,\y) {};
   \node[anchor=west,inner sep=2pt] at (\x,\y) {$\id$};
 }
 \node at (9.35,-.35) {$U_{a,p}\in[0,1)^2$, tag $(a,p)$};
\end{tikzpicture}
\caption{A schematic with $M=4$, $N=8$.  The column heights are the densities $Mn_a/N$ (height $1$ is uniform); the offsets retain the positions.  Together they recover the sample up to a permutation, via $x_{a,p}=J^{-1}(a+U_{a,p})$.  Each box on the right is a copy of the unit square; its points are numbered within that cell.  Under the uniform product reference, conditional on the counts these offsets are independent and uniform.}
\label{fig:occupancy-and-offsets}
\end{figure}


For this proof write $T=T_K$ and
$\ip{\sigma,T\tau}=\iint K(x,y)\,\dd\sigma(x)\dd\tau(y)$,
identifying a density with its measure when needed.  Continuity of $K$ makes this form finite on finite signed measures, and positivity extends from densities by approximation.  On the mean-zero cell-constant space $V_M$, define
\begin{equation}
 P_M(\nu_N)=\underset{h\in V_M}{\operatorname{argmin}}\,
                  \ip{h-\nu_N,T(h-\nu_N)}.
 \label{eq:explicit-histogram-square-shift}
\end{equation}
If the form has null directions, take the minimizer modulo them.  Cauchy--Schwarz for the nonnegative form shows that the mixed functional vanishes on these directions.  Projection orthogonality gives
\begin{equation}
 \ip{h,T(\nu_N-P_M(\nu_N))}=0
 \qquad(h\in V_M).
 \label{eq:histogram-square-shift}
\end{equation}
Put $z_N=f_L-1+P_M(\nu_N)$.  Expanding the energy and using \eqref{eq:histogram-square-shift} gives the completed square
\begin{equation}
 \begin{aligned}
 \ip{\mu_N-\muunif,T(\mu_N-\muunif)}
 &=\ip{z_N,Tz_N}+\ip{\nu_N,T\nu_N}
                   -\ip{P_M(\nu_N),TP_M(\nu_N)}\\
 &\le\ip{z_N,Tz_N}+\ip{\nu_N,T\nu_N}.
 \end{aligned}
 \label{eq:direct-square-completion}
\end{equation}

\textbf{2. Entropy bounds the completed square.}
For $h\in V_M$, the entropy variational formula and \eqref{eq:continuous-entropy-assumption} imply
\begin{equation}
 \begin{aligned}
 \log\int e^{Th}\,\dd\muunif
 &=\sup_\rho\left\{\ip{h,T(\rho-1)}-\Hr(\rho\muunif\mid\muunif)\right\}\\
 &\le\sup_\rho\left\{\ip{h,T(\rho-1)}
                         -\frac B2\ip{\rho-1,T(\rho-1)}\right\}
 \le\frac1{2B}\ip{h,Th}.
 \end{aligned}
 \label{eq:dual-log-hls}
\end{equation}
The last step completes the square in $\rho-1-h/B$.  Also, for every $h\in V_M$,
\begin{equation}
 \ip{h,Tz_N}=\int Th\,\dd\mu_N
            =\frac1N\sum_iTh(X_i),
 \label{eq:coarse-gaussian-pairing}
\end{equation}
by the definition of $\nu_N$ and projection orthogonality.

Choose coordinates on the nonnull part of $V_M$, of dimension $q_M\le M-1$, and let $\Sigma$ be the matrix of its energy form.  For a Gaussian coordinate vector $Z$ with covariance $\Sigma^{-1}$, let $\mathsf Z$ denote the corresponding cell-constant function.  The elementary identity
$\E_Z[e^{v^\top\Sigma Z}]=e^{v^\top\Sigma v/2}$ and \eqref{eq:coarse-gaussian-pairing} give, for $0\le t<B$,
\begin{equation}
 \begin{aligned}
 \E\big[e^{\frac{tN}2\ip{z_N,Tz_N}}\big]
 &=\E_Z\left[\left(\int e^{\sqrt{t/N}\,T\mathsf Z}
                           \,\dd\muunif\right)^N\right]\\
 &\le\E_Z\big[e^{\frac{t}{2B}\ip{\mathsf Z,T\mathsf Z}}\big]\\
 &=\frac1{(2\pi)^{q_M/2}\sqrt{\det(\Sigma^{-1})}}
   \int_{\R^{q_M}}e^{-\frac12(1-t/B)z^\top\Sigma z}\,\dd z\\
 &=(1-t/B)^{-q_M/2}.
 \end{aligned}
 \label{eq:coarse-partition-estimate}
\end{equation}
Here, the inequality is \eqref{eq:dual-log-hls} applied to
$h=\sqrt{t/N}\,\mathsf Z$.  Tonelli justifies exchanging the Gaussian and particle integrals.  The bound is independent of $N$ for the fixed grid.

\textbf{3. Continuity makes the remainder small.}
Set
\begin{equation}
 r_J=(I-Q_J)_x(I-Q_J)_yK.
 \label{eq:general-cell-centering}
\end{equation}
This subtracts the box average in each variable and adds back the joint average.  If $\omega_K(\varepsilon)$ is the modulus of continuity when both arguments move by at most $\varepsilon$, then
\begin{equation}
 \|r_J\|_\infty\le4\omega_K(\sqrt\d/J)\longrightarrow0.
 \label{eq:continuous-small-remainder}
\end{equation}
The kernel is centered in each variable under $\muunif$, and
$\ip{\nu_N,T\nu_N}=N^{-2}\sum_{i,j}r_J(X_i,X_j)$.
Removing equal particle indices in \eqref{eq:direct-square-completion}, we obtain
\begin{equation}
 \mathsf W_N(K)\le N\ip{z_N,Tz_N}
                  +\mathsf W_N(r_J)+5\|K\|_\infty.
 \label{eq:diagonal-correction-bound}
\end{equation}
Indeed, the diagonal correction is
$N^{-1}\sum_i[r_J(X_i,X_i)-K(X_i,X_i)]$, whose absolute value is at most $5\|K\|_\infty$.

Choose conjugate $p,q>1$ with $p\be_0<B$.  By Lemma~\ref{lem:canonical-array-exponential}, with $\kappa=\|r_J\|_\infty$, a sufficiently fine fixed grid gives
$\sup_{N,\,0\le\be\le\be_0}\E[e^{q\be\mathsf W_N(r_J)/2}]\le8/3$.
H\"older and \eqref{eq:coarse-partition-estimate} now yield
\begin{equation}
 \E\big[e^{\frac\be2\mathsf W_N(K)}\big]
 \le e^{5\be_0\|K\|_\infty/2}
      (1-p\be_0/B)^{-q_M/(2p)}(8/3)^{1/q}.
 \label{eq:continuous-proof-conclusion}
\end{equation}
This proves Lemma~\ref{lem:continuous-entropy-partition}.

For the singular kernel in Section~\ref{sec:local-singular}, the grid is instead tied to its support by $J^{-1}=2r$: only adjacent boxes interact, and the particle count makes their effective coupling small.

\begin{thebibliography}{1}
\bibitem{RenWang2025}
Z.~Ren and S.~Wang,
\emph{Size of chaos for Gibbs measures of mean field interacting diffusions},
Probability Theory and Related Fields (2025).
\href{https://doi.org/10.1007/s00440-025-01435-z}{doi:10.1007/s00440-025-01435-z}.
Accepted version: \href{https://arxiv.org/abs/2411.14236v3}{arXiv:2411.14236v3}.
\end{thebibliography}
\end{document}
